% 华侨大学项目建议书模板
% 编译方式：xelatex main.tex（建议连续编译两次以生成完整页眉/目录）
\documentclass{hqu-proposal}

% ================== 项目信息 ==================
\projectname{冷弯薄壁型钢有限元分析模块开发}
\applicant{吴俊超}
\college{土木工程学院}
\contact{13616008369 / jcwu@hqu.edu.cn}
\applydate{2026 年 7 月}
\logoleft{figures/logo1.pdf}
\logoright{figures/logo2.pdf}

% ================== 正文开始 ==================
\begin{document}

\makecover

\tableofcontents
\newpage

\section{项目基本信息}

\begin{table}[htbp]
    \centering
    \caption{项目基本信息表}
    \begin{tabular}{>{\heiti}p{3.5cm}p{4.5cm}>{\heiti}p{3cm}p{4cm}}
        \toprule
        项目名称 & \multicolumn{3}{l}{基于深度学习的校园智能安防系统研究与实现} \\
        \midrule
        项目编号 & 2026HQU-XXXX & 项目类型 & 创新训练项目 \\
        \midrule
        申请人 & 张三 & 学号 & 2026XXXXXX \\
        \midrule
        所在学院 & 计算机科学与技术学院 & 专业 & 计算机科学与技术 \\
        \midrule
        指导教师 & 李四 & 职称 & 教授 \\
        \midrule
        联系电话 & 138-xxxx-xxxx & 申请日期 & 2026 年 7 月 \\
        \bottomrule
    \end{tabular}
\end{table}

\section{项目简介}

\subsection{研究背景}
校园安全是高校管理工作的重中之重。随着高校规模的不断扩大，传统的人工巡逻和监控回放方式已难以满足实时预警、快速响应的需求。本项目拟结合深度学习技术，构建一套面向校园场景的智能安防系统，实现人员异常行为识别、重点区域入侵检测、人群密度估计等功能，提升校园安全管理水平。

\subsection{研究目的与意义}
\begin{enumerate}[label=(\arabic*)]
    \item \textbf{理论意义：}探索轻量级深度学习模型在边缘计算设备上的部署方法，丰富智能视频监控领域的研究。
    \item \textbf{应用价值：}为学校保卫部门提供实时、准确的预警信息，降低安全事故发生概率。
    \item \textbf{育人功能：}通过项目实践，培养学生的工程能力、创新意识和团队协作精神。
\end{enumerate}

\section{立项依据}

\subsection{国内外研究现状}
近年来，基于卷积神经网络（CNN）和 Transformer 的目标检测与行为识别算法发展迅速。YOLO 系列、Faster R-CNN 等模型在公开数据集上取得了优异性能；同时，基于视频理解的 SlowFast、TimeSformer 等网络在行为识别任务中表现突出。然而，这些模型在计算资源受限的校园监控设备上部署时仍面临实时性差、误报率高等问题。

\subsection{项目可行性分析}
\begin{itemize}[leftmargin=2em]
    \item \textbf{技术可行性：}团队具备 Python、PyTorch 和 OpenCV 开发经验，可完成模型训练与部署。
    \item \textbf{数据可行性：}可与学校保卫处合作，获取脱敏后的校园监控视频数据。
    \item \textbf{硬件可行性：}实验室配备高性能 GPU 服务器和 NVIDIA Jetson 边缘计算设备。
\end{itemize}

\section{研究方案}

\subsection{研究内容}
\begin{enumerate}[label=\arabic*., leftmargin=2em]
    \item 构建校园安防数据集，包括异常行为、非法入侵、人群聚集等类别。
    \item 设计并训练面向校园场景的轻量化检测模型。
    \item 开发实时视频分析平台，支持多路摄像头接入与告警推送。
    \item 在典型场景（校门、图书馆、宿舍区）开展系统测试与优化。
\end{enumerate}

\subsection{技术路线}
本项目的技术路线如图~\ref{fig:tech-route} 所示。首先对校园监控视频进行采集与预处理；其次基于 YOLOv8 与 LSTM 构建检测-跟踪-行为识别一体化模型；最后通过 Flask 后端与 Vue 前端实现可视化监控平台。

\begin{figure}[htbp]
    \centering
    \fbox{\parbox{0.8\textwidth}{\centering \vspace{2cm} 技术路线图（请替换为实际图片）\vspace{2cm}}}
    \caption{项目技术路线}
    \label{fig:tech-route}
\end{figure}

\subsection{研究方法}
本项目采用文献调研、实验验证、对比分析与工程实现相结合的研究方法。通过消融实验验证各模块有效性，并采用准确率（Accuracy）、精确率（Precision）、召回率（Recall）和实时帧率（FPS）作为评价指标。

\section{项目进度安排}

\begin{table}[htbp]
    \centering
    \caption{项目进度安排表}
    \begin{tabular}{clp{8cm}}
        \toprule
        \textbf{序号} & \textbf{时间节点} & \textbf{主要工作内容} \\
        \midrule
        1 & 2026.09--2026.10 & 文献调研、需求分析、数据集采集 \\
        2 & 2026.11--2026.12 & 模型设计、训练与调优 \\
        3 & 2027.01--2027.03 & 平台开发与典型场景部署 \\
        4 & 2027.04--2027.05 & 系统测试、性能优化与结题验收 \\
        \bottomrule
    \end{tabular}
\end{table}

\section{预期成果}

\begin{enumerate}[label=(\arabic*)]
    \item 完成一套校园智能安防原型系统，支持至少 4 路摄像头实时分析。
    \item 形成 1 份完整的技术报告和用户使用手册。
    \item 发表或投稿相关学术论文 1 篇（中文核心或 EI/CPCI 会议）。
    \item 申请软件著作权 1 项。
\end{enumerate}

\section{经费预算}

\begin{table}[htbp]
    \centering
    \caption{项目经费预算表}
    \begin{tabular}{clcr}
        \toprule
        \textbf{序号} & \textbf{支出科目} & \textbf{金额（元）} & \textbf{计算依据} \\
        \midrule
        1 & 设备费 & 3,000 & 边缘计算设备、存储硬盘等 \\
        2 & 材料费 & 1,000 & 实验耗材、办公用品等 \\
        3 & 测试化验加工费 & 500 & 模型训练云服务资源 \\
        4 & 差旅费 & 1,000 & 参加学术会议、调研交通 \\
        5 & 版面/知识产权费 & 1,500 & 论文版面费、软著申请费 \\
        \midrule
        \multicolumn{2}{c}{\textbf{合计}} & \textbf{7,000} & \\
        \bottomrule
    \end{tabular}
\end{table}

\section{项目团队}

\begin{table}[htbp]
    \centering
    \caption{项目团队信息表}
    \begin{tabular}{clllp{4cm}}
        \toprule
        \textbf{序号} & \textbf{姓名} & \textbf{学号} & \textbf{学院/专业} & \textbf{项目分工} \\
        \midrule
        1 & 张三 & 2026XXXXXX & 计算机学院/计科 & 项目负责人、算法设计 \\
        2 & 王五 & 2026XXXXXX & 计算机学院/计科 & 后端开发 \\
        3 & 赵六 & 2026XXXXXX & 计算机学院/软件 & 前端与可视化 \\
        4 & 孙七 & 2026XXXXXX & 计算机学院/人工智能 & 数据标注与测试 \\
        \bottomrule
    \end{tabular}
\end{table}

\section{参考文献}

\begin{enumerate}[label={[\arabic*]}, leftmargin=2em]
    \item Redmon J, Divvala S, Girshick R, et al. You only look once: Unified, real-time object detection[C]. CVPR, 2016: 779-788.
    \item Feichtenhofer C, Fan H, Malik J, et al. SlowFast networks for video recognition[C]. ICCV, 2019: 6202-6211.
    \item 刘明, 王强. 基于深度学习的视频监控异常行为检测综述[J]. 计算机学报, 2023, 46(5): 1123-1145.
\end{enumerate}

\end{document}
